\subsection{对偶的例子}

命题 17.　设 \(n \geqslant 1\) 是一个整数。用 \(\boldsymbol{\mu}_{n}\) 表示 \(n\) 次单位根（在 \(\mathbf{C}\) 中）所成的群。群 \(\mathbf{Z}/n\mathbf{Z}\) 与 \(\boldsymbol{\mu}_{n}\) 关于由映射 \(\mathbf{Z} \times \boldsymbol{\mu}_{n} \to \mathbf{U}\)（定义为 \((m,z) \mapsto z^{m}\)）过渡到商所诱导的映射成对偶。

群 \(\mathbf{Z}/n\mathbf{Z}\) 与同态 \(\chi\)（从 \(\mathbf{Z}/n\mathbf{Z}\) 到 \(\mathbf{U}\)）所成的集合重合。这些同态形如 \(m \mapsto \chi(1)^m\)，其中 \(\chi(1)\) 是 \(\mathbf{U}\) 中满足 \(\chi(1)^n = 1\) 的任意元素，由此即得结论。

推论 1.　设 G 是一个有限交换群。对偶群 \(\widehat{\mathbf{G}}\) 同构于 G。

群 G 同构于循环群的一个有限乘积（A，VII，第 22 页，定理 3），而它的对偶群同构于这些循环群的对偶群的乘积（II，第 206 页，命题 2）。于是归结为 G 是循环群的情形，而这由命题 17 得到，因为群 \(\mu_{n}\) 是 n 阶循环群（A，V，第 75 页，定理 1）。

推论 2.　设 G 是一个局部紧交换群。群 G 是有限的，当且仅当 \(\widehat{\mathbf{G}}\) 是有限的。G 的一个闭子群 H 具有有限指数，当且仅当它的正交是有限的。

由对偶性，第一个断言由有限群的对偶是有限的这一事实（推论 1）得到。第二个断言由它得到，因为 \(\widehat{\mathrm{G / H}}\) 与 \(\mathrm{H}^{\perp}\) 等同（II，第 226 页，定理 4）。

命题 18.　\(\widehat{\mathbf{G}}\) 为紧的必要充分条件是 \(\mathbf{G}\) 是离散的。若 \(\mathbf{G}\) 是离散的，则 \(\mathbf{G}\) 上计数测度的对偶测度是 \(\widehat{\mathbf{G}}\) 上的规范化 Haar 测度。若 \(\mathbf{G}\) 是紧的，则规范化 Haar 测度的对偶测度是 \(\widehat{\mathbf{G}}\) 上的计数测度。

假设 G 是离散的，设 \(\alpha\) 是 G 上的计数测度。设 \(\varphi\) 是 \(e \in G\) 的特征函数。我们有 \(\mathcal{F}_{\mathrm{G}}(\varphi) = 1\)（在 \(\widehat{G}\) 上）。由于 \(\mathcal{F}_{\mathrm{G}}(\varphi)\) 在无穷远处趋于 0，群 \(\widehat{G}\) 是紧的。

此外，对测度 \(\widehat{\alpha}\)（\(\alpha\) 的对偶测度）而言，函数 \(\mathcal{F}_{\mathrm{G}}(\varphi)\) 的积分应为 \(\varphi(e)=1\)（II，第 219 页，命题 12）。因此 \(\widehat{\alpha}(\widehat{\mathrm{G}})=1\)。

假设 G 是紧的。此时 Haar 测度 dx 属于 \(\mathcal{M}^{1}(G)\)。它的 Fourier 变换在 \(\chi = e\) 处严格为正，而对 \(\chi \neq 0\) 为零（II，第 210 页，命题 6）。由于它在 \(\widehat{G}\) 上连续，群 \(\widehat{G}\) 是离散的。若 G 的测度为 1，则由对偶性从上一情形推出，测度 dx 的对偶测度是 \(\widehat{G}\) 上的计数测度。

推论 1（正交关系）.　假设 G 是离散的且赋予计数测度。对 G 中的 x 与 y，我们有

\[
\int_ {\widehat {G}} \chi (x) \overline {{\chi (y)}} d \chi = \left\{ \begin{array}{l l} 0 & \text{若 } x \neq y \\ 1 & \text{若 } x = y. \end{array} \right.
\]

这由 II，第 215 页定理 1 的推论以及对偶性得到。

推论 2.　设 H 是 G 的一个闭子群。

\begin{enumerate}
\def\labelenumi{\alph{enumi})}
\item
  H 为紧的必要充分条件是 \(H^{\perp}\) 在 \(\widehat{G}\) 中是开的；
\item
  H 为开的必要充分条件是 \(H^{\perp}\) 在 \(\widehat{G}\) 中是紧的。
\item
  说 \(H^{\perp}\) 是开的，等于说 \(\widehat{G}/H^{\perp}\) 是离散的；而 \(\widehat{G}/H^{\perp}\) 同构于 \(\widehat{H}\)（II，第 226 页，定理 4），因此该断言由命题 18 得到。断言 b) 由把断言 a) 应用于 \(H^{\perp}\) 再经对偶性得到。
\end{enumerate}

推论 3.　设 \((\mathrm{H}_{i})_{i\in\mathrm{I}}\) 是 G 的紧子群的一个递降滤过族。G 与群 \(G/H_{i}\) 的射影极限等同的必要充分条件是 \(\hat{G}\) 为诸开子群 \(H_{i}^{\perp}\) 的并。

说 G 与 \(G/H_{i}\) 的射影极限等同，等于说 \(\bigcap_{i}H_{i}=\{e\}\)（TG，III，第 60 页，命题 2），也就是 \(\bigcup_{i}H_{i}^{\perp}\) 在 \(\widehat{G}\) 中稠密（II，第 226 页，定理 4 的推论 4）。而 \(\bigcup_{i}H_{i}^{\perp}\) 是 \(\widehat{G}\) 的一个开子群，从而是闭子群。

推论 4.　设 I 是一个集合，\((\mathrm{H}_{i})_{i\in\mathrm{I}}\) 是一族紧群。诸 \(H_{i}\) 的乘积群的对偶是诸群 \(\widehat{H}_{i}\) 的离散群直和。

这是推论 3 的一个特殊情形。

命题 19.　设 K 是一个局部紧、非离散的域，不必交换，并设 G 是 K 的加法群，其群运算写作加法。设 \(\chi\) 是 G 的一个异于 1 的酉特征标。对 \(x,y\in\mathbf{G}\)，令 \(\theta(x,y)=\chi(xy)\)。则 G 关于 \(\theta\) 与自身成对偶。

对 \(y \in \mathbf{G}\)，用 \(\chi_y\) 表示从 \(\mathbf{G}\) 到 \(\mathbf{U}\) 中且满足 \(\chi_y(x) = \chi(xy)\) 的映射。我们有 \(\chi_y \in \widehat{\mathbf{G}}\)，且须证明 \(\beta: y \mapsto \chi_y\) 是从 \(\mathbf{G}\) 到 \(\widehat{\mathbf{G}}\) 上的一个拓扑群同构。

映射 \(\beta\) 是从 G 到 \(\widehat{\mathbf{G}}\) 中的一个单射同态；它是连续的（TG，X，第 28 页，定理 3 应用于连续映射 \(\theta\)（从 \(\mathbf{G} \times \mathbf{G}\) 到 \(\mathbf{C}\) 中））。我们来证明 \(\theta\) 是到其像上的一个同胚。只需（II，第 200 页，引理 2）证明：对 K 中 0 的每个邻域 U，存在 \(e\) 在 \(\widehat{\mathbf{G}}\) 中的一个邻域 V，使得 \(\overline{\beta}^{-1}(\mathrm{V}) \subset \mathrm{U}\)。设 \(x \mapsto |x|\) 是 K 上定义 K 的拓扑的一个绝对值（AC，VI，§9，第 1 节，命题 1），并设 \(x_0 \in \mathrm{K}\) 使得 \(\chi(x_{0}) \neq 1\)；记 \(\eta = |\chi(x_{0}) - 1| > 0\)。设 U 是 K 中 0 的一个邻域。存在 \(\delta > 0\)，使得 U 包含由 \(y \in K\)（满足 \(|y| < \delta\)）所成的集合。设 V 是由特征标 \(\xi \in \widehat{G}\)（对它们有 \(|\langle\xi, x\rangle - 1| < \eta\)，只要元素 \(x \in K\) 满足 \(|x| \leqslant |x_{0}|/\delta\)）所成的集合。它是 \(\widehat{G}\) 中 e 的一个邻域。若 \(y \neq 0\) 使 \(\beta(y)\) 属于 V，则我们有 \(|\chi(xy) - 1| < |\chi(x_{0}) - 1|\)（对所有满足 \(|x| \leqslant |x_{0}|/\delta\) 的 x）。从而 \(|x_{0}y^{-1}| > |x_{0}|/\delta\)，于是 \(|y| < \delta\)，故 \(y \in U\)。

由于 \(\beta\) 是到其像上的一个同胚，它的像在 \(\widehat{\mathbf{G}}\) 中是闭的（TG，III，第 22 页，推论 2）。但另一方面，\(\beta\) 的像的正交是由元素 \(x\)（属于 \(\mathbf{G}\)，且 \(\chi(xy) = 1\) 对所有 \(y \in \mathbf{G}\) 成立）组成的集合，因而化为 \(\{0\}\)。于是 \(\beta\) 的像在 \(\widehat{\mathbf{G}}\) 中稠密（II，第 228 页，推论 3）。我们断定 \(\beta\) 是满射的。

推论 1.　设 K 是一个局部紧、非离散、不必交换的域，\(\chi\) 是 K 的加法群的一个非平凡酉特征标。设 E 是 K 上一个有限维拓扑向量空间。映射 \(\theta\)（从 \(E \times E'\) 到 U 中，由 \(\theta(x,\lambda)=\chi(\langle\lambda,x\rangle)\)（对 \((\lambda,x)\in\mathrm{E}'\times\mathrm{E}\)）定义）使拓扑群 E 与 E' 成对偶。

设 n 是 E 的维数，\((e_{1},\ldots,e_{n})\) 是 E 的一个基。它使我们可以把 E 与 \(K^{n}\) 等同（EVT，I，第 14 页，定理 2）。于是结论由命题 19 与 II，第 206 页命题 2 得到。

我们用 T 表示群 R/Z。

推论 2.　a) 群 R 关于映射 \((x,y)\mapsto\exp(2i\pi xy)\) 与自身成对偶，且 Lebesgue 测度的对偶测度是 Lebesgue 测度；

\begin{enumerate}
\def\labelenumi{\alph{enumi})}
\setcounter{enumi}{1}
\tightlist
\item
  群 Z 与 T 关于由从 \(Z \times R\) 到 U 且 \((n, x) \mapsto \exp(2i\pi nx)\) 的映射过渡到商所得到的映射成对偶。Z 上计数测度的对偶 Haar 测度是 R/Z 上的规范化 Haar 测度。
\end{enumerate}

由命题 19，群 \(\mathbf{R}\) 关于映射 \((x,y)\mapsto \exp (2i\pi xy)\) 与自身成对偶。我们把 \(\widehat{\mathbf{R}}\) 与 \(\mathbf{R}\) 等同。\(\mathbf{Z}\) 在 \(\widehat{\mathbf{R}} = \mathbf{R}\) 中的正交就是 \(\mathbf{Z}\)，而 \(b)\) 由 II，第 226 页定理 4 得到。

设 \(\alpha\) 是 Z 上的计数测度，\(\gamma\) 是 T 上的规范化 Haar 测度。若用 \(\beta\) 表示 \(\mathbf{R}\) 上的 Lebesgue 测度，我们有 \(\gamma = \beta/\alpha\)，因为 R/Z 上这两个 Haar 测度的质量都是 1。Haar 测度 \(\widehat{\alpha}\)（在 \(\widehat{Z} = T\) 上）是规范化 Haar 测度（命题 18），而 Haar 测度 \(\widehat{\gamma}\) 是 \(\mathbf{Z}\) 上的计数测度（同前）。由 II，第 231 页命题 16，\(\beta\) 的对偶测度因此是测度 \(\beta\)。

注.　特别地，我们重新得到 I，第 36 页例 4 中对 \(\mathsf{X}(\mathsf{L}^{1}(\mathbf{Z}))\) 的确定。

对每个整数 \(n \geqslant 0\) 及 \((x, y) \in \mathbf{R}^n \times \mathbf{R}^n\)，我们记

\[
x \cdot y = \sum_ {j = 1} ^ {n} x _ {j} y _ {j}.
\]

推论 3.　设 \(n \geqslant 1\) 是一个整数。群 \(\mathbf{R}^{n}\) 关于映射 \((x,y) \mapsto \exp(2i\pi x \cdot y)\) 与自身成对偶，且 \(\mathbf{R}^{n}\) 上 Lebesgue 测度的对偶测度是 Lebesgue 测度。群 \(\mathbf{Z}^{n}\) 与 \(\mathbf{T}^{n} = \mathbf{R}^{n}/\mathbf{Z}^{n}\) 关于由映射 \((n,x) \mapsto \exp(2i\pi x \cdot y)\) 过渡到商所得到的映射成对偶，且 \(\mathbf{Z}^{n}\) 上计数测度的对偶 Haar 测度是 \((\mathbf{R}/\mathbf{Z})^{n}\) 上的规范化 Haar 测度。

这由 II，第 225 页引理 7、II，第 217 页命题 10 以及推论 2 得到。

注.　对 \(R^{n}\) 的一个子群 H，相应地有它的正交 \(H^{\perp}\)，即 \(\widehat{R^{n}} = R^{n}\) 的一个子群，它正是 TG，VII，第 6 页，第 3 节中定义的与 H 相伴的子群。

在下文中，我们将借助该推论的对偶，把 \(R^{n}\)（相应地，\(T^{n}\)）的对偶与 \(R^{n}\)（相应地，\(Z^{n}\)）等同。特别地，对 \(f \in \mathrm{L}^{1}(\mathbf{R}^{n})\)，它的 Fourier 变换等同于 \(R^{n}\) 上取值于 C 的、把 \(y \in R^{n}\) 映为为

\[
\mathcal {F} (f) (y) = \int_ {\mathbf {R} ^ {n}} f (x) \exp (- 2 i \pi x \cdot y) d x.
\]

推论 4.　群 \(\mathbf{R}^{*}\) 与群 \(\{-1,1\} \times \mathbf{R}\) 通过映射 \((x,(\sigma,t)) \mapsto \sigma(x/|x|)|x|^{it}\) 成对偶。群 \(\mathbf{R}_{+}^{*}\) 与 \(\mathbf{R}\) 通过映射 \((x,t) \mapsto x^{it}\) 成对偶。

事实上，映射 \(x \mapsto (x/|x|, \log(|x|))\) 是从 \(R^{*}\) 到 \(\{-1, 1\} \times R\) 上的一个拓扑群同构。于是该断言由 II，第 225 页引理 7、推论 2 以及 \(\{-1, 1\}\) 的酉特征标恰为 1 与 \(x \mapsto x\) 这一事实得到。

设 \(p\) 是一个素数。域 \(\mathbf{Q}_{p}\)（即 \(p\) 进数域）是 \(\mathbf{Q}\) 关于 \(p\) 进赋值的完备化（INT，VII，§1，第 6 节，例，及 AC，VI，§3，第 4 节，例 4）。对每个 \(x \in \mathbf{Q}_{p}\)，存在唯一的整数 \(\nu\geqslant0\) 与唯一的满足 \(0\leqslant q<p^{\nu}\) 的整数 q，使得 \(qp^{-\nu}-x\in Z_{p}\)（A，VII，第 10 页，定理 2，应用于主理想整环 \(Z_{p}\) 与集合 \(R_{p}\)（由满足 \(0\leqslant j<p\) 的整数 j 组成））。我们记 \(\lambda(x)=qp^{-\nu}\)。

命题 20.　映射 \(x \mapsto \exp(2i\pi\lambda(x))\) 是 \(\mathbf{Q}_{p}\) 的一个酉特征标，其核为 \(\mathbf{Z}_{p}\)。

对 \(x_1\) 与 \(x_2\)（属于 \(\mathbf{Q}_p\)），按定义有 \(\lambda(x_1 + x_2) - \lambda(x_1) - \lambda(x_2) \in \mathbf{Z}_p \cap \mathbf{Q} = \mathbf{Z}\)。此外，映射 \(\lambda\) 是局部常数的，因为 \(\lambda(x + y) = \lambda(x)\)（若 \(y \in \mathbf{Z}_p\)）。由此推出 \(x \mapsto \exp(2i\pi\lambda(x))\) 是 \(\mathbf{Q}_p\) 的一个酉特征标。由于 \(\lambda(x) \in \mathbf{Z}\) 当且仅当 \(x \in \mathbf{Z}_p\)，这个特征标的核是 \(\mathbf{Z}_p\)。

回顾：\(Q_{p}\) 的加法群上的规范化 Haar 测度是指唯一的 Haar 测度 \(\mu\)，它使 \(\mu(\mathbf{Z}_{p}) = 1\)（INT，VII，§1，第 6 节，例）。

推论.　a) 群 \(\mathbf{Q}_{p}\) 关于映射 \((x,y)\mapsto\exp(2i\pi\lambda(xy))\) 与自身成对偶。\(\mathbf{Q}_{p}\) 上的规范化 Haar 测度这时是它自己的对偶；

\begin{enumerate}
\def\labelenumi{\alph{enumi})}
\setcounter{enumi}{1}
\tightlist
\item
  群 \(Z_{p}\) 与 \(Q_{p}/Z_{p}\) 关于由 \((z,x)\mapsto\exp(2i\pi\lambda(zx))\) 所定义的映射过渡到商而得到的映射成对偶，且 \(Z_{p}\) 上规范化 Haar 测度的对偶测度是 \(Q_{p}/Z_{p}\) 上的计数测度。
\end{enumerate}

其证明逐字重复命题 19 的推论 2 的证明。

